% ============================================================
% 01-introduction.tex
% ============================================================
\section{Introduction}
\label{sec:intro}

Providing timely, accessible, and accurate administrative guidance is vital for student success and institutional efficiency. In large universities such as Can Tho University (CTU)---a flagship institution in Vietnam's Mekong Delta serving over 40,000 students~\cite{ctu2022handbook}---advisory inquiries span diverse administrative domains. Because regulations are frequently updated through ministerial circulars, rector decisions, and cohort-specific guidelines, automated conversational counseling systems are essential to alleviate advisory strain and ensure compliance.

However, higher education counseling is characterized by acute \emph{knowledge heterogeneity}, encompassing structurally distinct representations: (1)~\textbf{Academic Curricula:} relational course prerequisite chains spanning the 113 majors represented in our graph corpus; (2)~\textbf{Tuition Schedules:} multi-dimensional rate matrices governed by admission cohort and discipline; (3)~\textbf{Statutory Exemptions:} legal decrees prescribing percentage deductions for qualified beneficiaries~\cite{government2021decree81,government2023decree97}; (4)~\textbf{Scholarships:} multi-variable eligibility thresholds (cumulative GPA and training point scores); and (5)~\textbf{Campus Regulations:} narrative administrative prose governing examinations, conduct, and student loans.

To automate institutional QA, Retrieval-Augmented Generation (RAG)~\cite{lewis2020rag,gao2023rag_survey} and GraphRAG~\cite{edge2024graphrag} augment large language models (LLMs) with external documents or knowledge graphs, while agentic architectures~\cite{yao2023react,wu2023autogen,schick2023toolformer} introduce iterative tool invocation and multi-agent coordination. Nevertheless, applying existing paradigms to university counseling reveals significant limitations. Monolithic RAG pipelines risk a \emph{uniform representation fallacy}: forcing relational curricula, tabular fee matrices, and narrative legal prose into a single retrieval index can fragment tabular data and prerequisite links. Conversely, uncoordinated multi-agent meshes can expose broad tool sets across agents, increasing the risk of tool mis-selection, parameter errors, and extended execution, while known limitations in language-model arithmetic~\cite{qian2023limitations} motivate deterministic calculation for credit-based tuition.

These challenges crystallize three open research gaps in automated university counseling:
(1)~\textbf{Gap 1 (Limited Domain-Specialized Multi-Agent Orchestration):} Existing agentic and multi-agent systems provide delegation and tool use, but the coordinated use of domain-specialized agents with isolated knowledge spaces, tool privileges, and retrieval paths remains underexplored in heterogeneous university counseling;
(2)~\textbf{Gap 2 (Uniform Knowledge Representation):} Existing RAG-oriented systems often process structurally different institutional knowledge through a common retrieval space, although relational curricula, structured tuition information, and narrative regulations require fundamentally different representations; and
(3)~\textbf{Gap 3 (Lack of Representation-Aware Routing):} Existing systems provide lexical, semantic, and graph retrieval methods, but offer limited mechanisms for selecting the appropriate evidence representation and specialist according to query domain and information structure.

To address these gaps, we formulate three core research questions:
\begin{itemize}
    \setlength{\itemsep}{1.5pt}
    \setlength{\parsep}{0pt}
    \setlength{\parskip}{0pt}
    \item \textbf{RQ1 (Domain-Specialized Multi-Agent Architecture):} How effectively can a centralized supervisor-routed multi-agent architecture coordinate domain-specialized agents with isolated knowledge and tool access for heterogeneous university counseling tasks?
    \item \textbf{RQ2 (Heterogeneous Knowledge Allocation):} How does prioritizing graph-oriented access for relational academic-program and structured tuition knowledge, while retaining narrative regulations in hybrid document retrieval, affect retrieval and answer quality compared with uniform retrieval representations?
    \item \textbf{RQ3 (Representation-Aware Evidence Routing):} How does intent-guided routing between graph-oriented and hybrid document-retrieval paths affect retrieval effectiveness, answer correctness, and latency across different counseling domains?
\end{itemize}

\emph{We hypothesize that reliable university counseling is better supported by a supervisor-routed multi-agent architecture in which domain specialists are constrained to knowledge representations and tools appropriate to their information structure, rather than by a uniform retrieval-and-reasoning pipeline. Accordingly, \system{} separates specialist coordination, knowledge allocation, and evidence-path selection as three complementary design dimensions.}

To answer these questions, we present \textbf{\system{}}. \emph{\system{} proposes a centralized supervisor-routed, domain-specialized multi-agent architecture in which each specialist is assigned a domain-specific knowledge space, retrieval path, and tool privilege.} Under centralized supervisor coordination, incoming inquiries are contextualized via Redis history and dispatched to four bounded specialists operating under asymmetric privileges: (1)~\texttt{Academic Specialist}, bound to 6 dedicated Cypher tools over a Neo4j knowledge graph; (2)~\texttt{Financial Specialist}, equipped with structured tuition graph lookups and the deterministic arithmetic engine \texttt{tinh\_toan\_hoc\_phi}; (3)~\texttt{Scholarship Specialist}, equipped with the deterministic evaluator \texttt{tinh\_tien\_hoc\_bong}; and (4)~\texttt{General Specialist}, operating with zero tools to synthesize narrative regulatory context. Specialists share session state in Redis while maintaining strictly isolated execution envelopes without unconstrained peer-to-peer communication.

Our principal contributions are three-fold:
\begin{itemize}
    \setlength{\itemsep}{1.5pt}
    \setlength{\parsep}{0pt}
    \setlength{\parskip}{0pt}
    \item \textbf{C1 (Supervisor-Routed Domain-Specialized Multi-Agent Architecture):} A centralized supervisor-routed architecture coordinating four domain-specialized agents with asymmetric tool privileges, a per-decision bound of $\text{max\_tool\_calls}=1$, and shared session state, designed to reduce tool mis-selection and constrain execution without peer-to-peer delegation (addressing Gap~1 and RQ1).
    \item \textbf{C2 (Heterogeneous Knowledge Allocation Framework):} A principled allocation strategy that separates relational curricula and structured tuition into graph-oriented representations while reserving hybrid document retrieval for narrative administrative decrees, resolving the uniform representation fallacy (addressing Gap~2 and RQ2).
    \item \textbf{C3 (Representation-Aware Routing and Evidence Acquisition):} An intent-guided evidence acquisition mechanism that routes queries dynamically to graph-oriented or hybrid document-retrieval paths, leveraging established BM25, Dense vectors, RRF, and cross-encoder reranking as complementary building blocks (addressing Gap~3 and RQ3).
\end{itemize}
